\section{Simulated line, soft sensor and optimizer}\label{sec:methods}

\subsection{Simulated brine line}\label{sec:plant}

An estimator or an optimizer tested against a plant that behaves
exactly as its model assumes proves little. The simulated line
therefore uses the kinetics of Section~\ref{sec:kinetics} with
structured deviations, each parameterized and documented. The true
polymerization rate is
\begin{equation}\label{eq:truth}
r_{\mathrm{true}} \;=\; f_u(u)\, e^{B(T)}\, k(T)\, s^{\,n},
\qquad s = C - C_e(T)\,\beta_e ,
\end{equation}
gated by an induction clock. Here $k(T)$ and $C_e(T)$ are the
literature-fitted rate constant and solubility; $B(T)$ is a smooth
log-normal bias drawn per plant on the seven kinetic isotherm nodes
($\sigma = 0.35$); the reaction order is drawn as
$n \sim \mathcal{N}(4, 0.15)$ and clipped to $[3.6, 4.4]$, in keeping
with the ``approximately fourth order'' of the source data;
$\beta_e \sim \mathcal{N}(1, 0.03)$ perturbs the solubility; and $f_u$
is a per-plant instance of the non-monotonic dose--response family of
Section~\ref{sec:kinetics}, with the maximum suppression drawn about
0.50 (standard deviation 0.10), the dose of maximum effect drawn
log-normally about 1~ppm (log standard deviation 0.25), and a mild
reversal on overdose. The induction period is fitted from the primary
kinetic dataset, $\log_{10} P = a + b \log_{10} s + cT$ with
$b = -4.31$ and $r^2 = 0.88$, a relation that reproduces the source
publication's own graphical correlation. The plant keeps a
supersaturation-exposure clock that the model lacks.

The mixture weights of the four endmembers follow a mean-reverting
Ornstein--Uhlenbeck process about the operating composition, and the
flash fraction fluctuates within $[0.20, 0.30]$, so that the separated
brine wanders inside the flashed envelope while remaining stationary
(a drift of about 1.3\% over ten days). An earlier random-walk form,
which drifted toward the mixture centroid and diluted the brine over
time, was replaced, and stationarity is now one of the repository's
automated checks.

The measurement layer follows what a shift-based operation provides:
silica grab samples every 8~h at 5\% relative noise, a continuous line
temperature with $\sigma = 0.3$~$^{\circ}$C, and one daily measurement
of the trailing-24-h mean polymerization rate at 10\% relative noise.
The rate measurement stands for a daily polymerization test, in which
the decline of monomeric silica in a brine sample is followed at line
conditions as in field inhibitor trials~\citep{baltazar2014grc}, or for
an online scale sensor~\citep{okazaki2017fiber}. A fourteen-day episode
thus provides 42 concentration samples and 14 rate labels against
2{,}016 hidden values of the true rate. A step-wise interface reproduces
the batch simulator's trajectory to machine precision, so that in
closed loop the optimizer decides each day from what it has observed.

\subsection{Hybrid soft sensor}\label{sec:sensor}

The sensed quantity is the trailing-24-h mean polymerization rate, the
quantity that accumulates scale and the one the optimizer constrains.
The sensor is a parallel hybrid~\citep{thompson1994modeling,%
vonstosch2014hybrid},
\begin{equation}\label{eq:sensor}
\hat r \;=\; \max\!\Big[\,0,\;
\exp\!\big(\log(1{+}r_{\mathrm{core}}) + \lambda\, g_{\theta}(x)\big) - 1
\Big],
\end{equation}
in which $r_{\mathrm{core}}$ is the trailing mean of the kinetic core
combined with the nominal dose--response, and $g_{\theta}$ is a
gradient-boosted residual~\citep{ke2017lightgbm} in $\log(1{+}\cdot)$
space over causal trailing-window aggregates of the measurements only;
a code-level guard raises an error if any hidden plant variable reaches
the feature set. Each sensor is trained on its own plant's sparse daily
labels under strict temporal splits, the standard soft-sensor
protocol~\citep{kadlec2009data}.

The protocol matters. Asked to generalize zero-shot to unseen plant
realizations, the residual lost to both a pure-physics and a pure-data
baseline. The reason is structural: per-plant mismatch cannot be
predicted before any data from that plant exist, whereas a sensor
trained on its plant's history and predicting forward in time can
identify that plant's bias. Two protective mechanisms follow. The
residual weight $\lambda$ in Eq.~\eqref{eq:sensor} is set by
leave-one-out (LOO) credibility on the training labels,
\begin{equation}\label{eq:lambda}
\lambda^{\ast} = \min\!\Big(1,\, \max\!\Big(0,\,
\frac{\langle e, \tilde r\rangle}{\langle \tilde r, \tilde r\rangle}\Big)\Big),
\quad \text{with } \lambda = 0 \text{ unless the LOO SSE improves by at least } 15\%,
\end{equation}
where $e$ collects the residual targets and $\tilde r$ their LOO
predictions, so that the correction is used only to the extent that it
has shown out-of-sample signal and the sensor falls back to its physics
prior when labels are few or uninformative. The same LOO computation
gives the residual's prediction-error standard deviation
$\sigma_{\mathrm{LOO}}$, which the optimizer can use as the sensor's
own statement of uncertainty.

\subsection{Daily optimization of temperature and dose}\label{sec:rto}

The decision variables are the reinjection temperature set-point and
the inhibitor dose, $(T_r, u)$, refreshed once a day following the
two-layer structure of steady-state real-time
optimization~\citep{darby2011rto}. The objective is the revenue of the
power in Eq.~\eqref{eq:power} minus the inhibitor cost. The scaling
constraint is enforced on every vertex of the brine envelope at the
adverse flash bound,
\begin{equation}\label{eq:rto}
\max_{T_r,\,u}\; J(T_r, u)
\quad \text{s.t.} \quad
r_v(T_r, u) \;\le\; r_{\max} \quad \forall v \in \mathcal{V},
\end{equation}
where $r_v = k(T_r)\,(C_v - C_e(T_r))^4 f(u)$ is the kinetic core at
the silica concentration $C_v$ of vertex $v$, $f$ is the nominal
dose--response, and $r_{\max} = 0.5$~g\,t$^{-1}$\,min$^{-1}$
(Section~\ref{sec:units} relates it to wall growth). The problem is
solved by dense grid search over 95--175~$^{\circ}$C and 0--4~ppm
($161 \times 81$ points, steps of 0.5~K and 0.05~ppm): on a bounded
two-dimensional decision space the grid is global, deterministic and
auditable, and its vectorized form is checked against the reference
solver at start-up. Solved once and held, Eq.~\eqref{eq:rto} is the
open-loop robust policy.

\subsection{Feedback corrections compared}\label{sec:feedback}

Once a day the plant returns one noisy label $\ell_t$ of the
trailing-24-h mean rate. Three ways of using it are compared, all of
which correct the constraint of Eq.~\eqref{eq:rto} and leave the
objective unchanged. The gradient corrections of full modifier
adaptation~\citep{marchetti2009modifier,marchetti2010dual} are not
attempted, because one noisy label per day does not support estimating
plant gradients.

\subsubsection{Scalar correction and its anchor}\label{sec:anchor}

The zeroth-order scheme of \citet{chachuat2008constraint} filters the
constraint bias,
\begin{equation}\label{eq:ca}
\varepsilon_t = (1-\alpha)\,\varepsilon_{t-1}
 + \alpha\,\big(\ell_t - r_a(T_{r,t}, u_t)\big)
\quad \text{and} \quad
r_v + \varepsilon_t \le r_{\max},
\end{equation}
with $\alpha = 0.5$. A log-space variant filters
$\ln \ell_t - \ln r_a$ and scales the constraint by $e^{\theta_t}$. The
model rate $r_a$ against which the label is compared is the
\emph{anchor} of the correction, and the results turn on it. With the
\emph{design-case anchor}, $r_a = r_B$: the kinetic core at the
composition the optimizer is solved against, the envelope vertex
nearest the operating brine at the adverse flash bound. With the
\emph{measured anchor}, $r_a = k(T_r)\,(C_m - C_e(T_r))^4 f(u)$: the
core at the latest silica grab sample $C_m$.

\subsubsection{Gaussian-process residual}\label{sec:gp}

The constraint model becomes $\log(1+\hat r_v) = \log(1+r_v) + \mu(x)$,
where $\mu$ is the posterior mean of a Gaussian process trained on the
residual targets $\log(1+\ell_t) - \log(1+r_a)$ at the measured anchor,
with inputs $x = (T_r, u, C_m - C_e)$, a constant mean and a
squared-exponential kernel with one length scale per input plus noise.
Hyperparameters are chosen by exact marginal likelihood over a fixed
grid, so the fit is deterministic. This arm takes from modifier
adaptation with Gaussian
processes~\citep{ferreira2018gpma,delriochanona2021tale} the idea of a
learned mismatch model that reports its own uncertainty. It does not
implement those methods' cost modifiers, acquisition functions or trust
regions, and it tests the regressor, not those algorithms.

\subsubsection{Hybrid sensor}\label{sec:hybridarm}

The trained sensor of Section~\ref{sec:sensor} is used as the
constraint model, $\log(1+\hat r_v) = \log(1+r_v) + \lambda\,
g_\theta(x)$, queried at steady state across the decision grid. Its
residual target and anchor are those of the Gaussian-process arm. The
two differ in the regressor and in the inputs: the hybrid receives the
sensor's full set of trailing-window features, the Gaussian process the
three listed above.

\subsubsection{Uncertainty back-off}\label{sec:backoff}

Each correction can be backed off by a gain $\kappa$ times its own
uncertainty estimate,
\begin{equation}\label{eq:backoff}
\text{scalar: } r_v + \varepsilon_t + \kappa\,\sigma_u \le r_{\max},
\qquad
\text{learned: } \log(1+\hat r_v) + \kappa\,\sigma \le \log(1+r_{\max}),
\end{equation}
where $\sigma_u$ is the standard deviation of the scalar scheme's
one-step prediction error over the trailing 14 days, and $\sigma$ is the
Gaussian-process predictive standard deviation or the hybrid's
$\sigma_{\mathrm{LOO}}$. The gain is swept over
$\kappa \in \{0, 0.5, 1, 2\}$. The three uncertainty estimates have
different scales, so equal $\kappa$ does not mean equal caution, and the
families are compared as curves over the sweep, not gain by gain. The
learned models are refitted weekly once seven labels exist. Until then
they run the measured-anchored scalar correction restricted to
tightening ($\varepsilon \ge 0$), so that feedback cannot relax the
constraint before a model has earned credibility.

\subsection{Simulation experiments}\label{sec:design}

The estimation experiment simulates \NUM{T1plants} plant realizations
for 90 days each; sensors are trained on the labels of the first 60
days and tested on the final 30.

The closed-loop experiment runs every policy on \NUM{CLplantsB} plant
realizations for \NUM{CLdays} days in each of two regimes. The
\emph{benign} regime uses the random per-plant deviations of
Section~\ref{sec:plant}. The \emph{stress} regime is a single-factor
perturbation: the plant's true rate is six times the model's, with all
random draws set to zero. Its size is measured, not chosen; it is the
median gap of Section~\ref{sec:labtest}. A violation day is a day whose
realized trailing-24-h mean rate exceeds $r_{\max}$. Violations are
counted on the realized plant, not on the worst-case vertex of the
envelope, from day 2 and, after the 8-day warm-up of the learned models,
from day 9. Power gains are relative to fixed practice. Every policy
meets the same brine and mismatch realization on a given plant, so
policies are compared by paired differences over plants, with seeded
percentile-bootstrap 95\% intervals (4{,}000 resamples, plants as the
resampling unit).

Two kinds of evidence are kept apart. On six of the twenty plants per
regime, the hybrid arm at $\kappa \in \{0, 1\}$ and the open-loop
policy formed a protocol whose pass criteria were fixed before it was
run; its verdicts are reported as registered in
Section~\ref{sec:models}. Everything else in
Sections~\ref{sec:anchoring} to~\ref{sec:operating}, namely the
anchoring ablation, the scalar and Gaussian-process families, the gain
sweep and the fourteen further plants, was added afterwards with
earlier results in view, and is post hoc.
